% PATCH: C-07
% Target section: sections/results.tex (intro warehouse-sim sentence, STL-discretization
%   paragraph, appetitive real-vs-sim paragraphs, tab:consolidated_physical_metrics,
%   tab:terrain_sequence, all "Differential Rolling"/"Differential Mobile" instances)
% Requirement: intake/pending/R02-01-06_limpieza_editorial_menor.md, §1 (residual
%   text<->table/figure redundancy) + §2 (grammar/style), executed in full -- low priority, no
%   blockers, independent of every other R02-01 sub-block.
%
% §1.1 -- tab:consolidated_physical_metrics, Complex row, "Var(z) peak" column:
%   Before: "Three-phase, peak $>30{,}000$" (no unit, no reference scale against the other 3
%     rows: ~1100, dual-peak with no number, sustained).
%   After:  \revblue{"Three-phase, peak ${\approx}\,27\times$ Obstacle baseline"} -- author
%     confirmed the relative-scale option (over a footnote-with-unit or leaving it untouched);
%     ~30,000 / ~1100 (Obstacle row) ~= 27.
%
% §1.2 -- FigReal_Terrain_Latencies_4Panels.png, 3-paragraph redundancy (originally flagged
%   against results.tex:945-949 in the now-superseded audit): re-checked against the current
%   file -- only 2 paragraphs remain after the figure (an earlier pass had already condensed the
%   purely-descriptive setup paragraph into the interpretive one). No change needed; noted as
%   already resolved rather than re-merging non-existent redundant text.
%
% §2.1 -- results.tex:9 (word-level, no \revblue{} -- zero meaning change):
%   Before: "its scope and caveats as a symbolic, non-certified test arena are discussed in..."
%   After:  "its scope and caveats as a symbolic representation rather than a certified test
%     arena are discussed in..."
%
% §2.2 -- STL-mesh discretization paragraph (structural move, \revblue{}): was stranded after
%   the second stability test's neural-response paragraph, between prose and a figure block,
%   disconnected from the inclined-slope test setup it actually describes. Relocated verbatim to
%   the end of the paragraph that introduces that setup ("In a second test, the transition from
%   flat terrain to an inclined slope is defined..."), removed from its old position. No wording
%   changed, only location -- wrapped in \revblue{} per the project convention for structural
%   reorganization.
%
% §2.3 -- appetitive real-vs-simulation paragraphs (structural fusion, \revblue{}): two
%   consecutive paragraphs used near-identical contrastive phrasing ("unlike the simulation,
%   where..." / "In contrast, in the simulation...") to make two different points (behavioural
%   timing vs. neural-activation detail). Merged into one paragraph preserving both facts (the
%   23-second timing + neuron $X_5$/$X_7$/$X_{11}$ activations for the physical robot; neuron
%   $X_{12}$'s early activity for the simulation) without the repeated "unlike/in contrast"
%   template.
%
% §2.4 -- Limitations run-on sentence (originally flagged as >150 words with chained "yet",
%   "rather than", "Instead" clauses): re-checked against the current file -- the paragraph is
%   already split into 6 sentences of reasonable length by an earlier pass (predates this patch).
%   No change needed; noted as already resolved rather than re-splitting an already-short
%   paragraph.
%
% §2.5 -- mode-name capitalization consistency (author-confirmed decisions, both beyond the
%   original item's scope of "just capitalization"):
%   1. Naming inconsistency found while surveying every mode-name occurrence: "Differential
%      Rolling" (introduction/methodology/simulation prose, and the majority of results.tex) vs.
%      "Differential Mobile" (physical/hardware subsection, including tab:mode_steady_state_energy
%      and the FigReal_Morphological_Transition_Energy caption/paragraph added in N-05) were two
%      distinct names for the same locomotion mode, not a case variant of one name -- one caption
%      (line ~737, "differential rolling (mobile) mode") had already hedged between the two.
%      Author confirmed: unify to "Differential Rolling" (the name used in the majority of the
%      document). All "Differential Mobile"/"differential mobile" instances renamed; the
%      now-redundant "(mobile)" parenthetical in the hedged caption removed.
%   2. Title Case applied consistently for all three mode names used as proper nouns (naming a
%      mode as a category -- enumerations, transition diagrams "X -> Y", table "Mode" columns):
%      "Quadrupedal Walking", "Differential Rolling", "Omnidirectional". Generic adjectival usage
%      in flowing prose ("quadrupedal mode", "differential rolling mode", "omnidirectional gait")
%      was left lowercase, unchanged -- this was already consistent throughout and not part of
%      the reported inconsistency. Concretely re-cased: tab:terrain_sequence's "Mode" column (6
%      rows, previously lowercase "Quadrupedal walking"/"Differential rolling"/"Rolling ->
%      stagnation") to Title Case, matching the Title Case already used in the gait-transition
%      figures/captions elsewhere in the same section.
%
% Verification: scripts/validate_tex.sh and scripts/check_roundtrip.sh both pass after
%   reassemble.py. No LaTeX toolchain available locally to run compile_pdf.sh.
% Status: applied
